\section{Recorded Inertial Sensor Protocol and Results}
\label{app:imu}

\paragraph{Data and labels.}
BROAD contains 39 trials recorded at 285.7~Hz with gyroscope, accelerometer,
and magnetometer readings and an optical reference orientation
\citep{laidig2021broad}. We use the reference only for labels. The heading is
the angle of the sensor's $x$ axis projected onto the horizontal plane. Steps
with missing optical data or pitch above $75^\circ$ are excluded from the loss
and the metrics but remain in the inputs. Averaging the IMU over blocks of ten
samples gives 28.6~Hz. The test split contains trials 02, 05, 07, 09, 11,
14, 16, 18, 20, 23, 25, 27, 29, 31, 33, 35, 37, and 39. Trials 03, 12, and 21
serve only for selecting the learning rate and model variant. The remaining
18 trials form the training split. Trials 28--38 contain magnetic disturbances,
and the split places 29, 31, 33, 35, and 37 in the test set.

\paragraph{Inputs.}
A complementary filter with gain $0.05$, tuned on training trials, runs
causally over each whole trial \citep{mahony2008complementary}. Its gravity
direction and the wrapped change $u_t$ of its heading join the gyroscope and
accelerometer features. Its absolute heading is not used. Task~B adds the
tilt-compensated magnetometer heading, the field norm, and the dip angle. The
frame offset between the magnetometer and optical headings is estimated on
undisturbed training trials only.

\paragraph{Models and training.}
Every learned model applies the same two-layer encoder of width 32 to the
standardized features. The input-set rotation uses $16$ complex channels with
angles $n\,u_t$ for $n=1,\dots,8$, two channels per harmonic, a learned initial
phase, and a Born readout over 16 sectors. Its readout and initial phase have
528 parameters. The encoder receives no gradient in this variant. The
learned-angle variant sets the angles to a linear function of the encoder output
initialized at the harmonic solution (2{,}480 parameters). A third variant adds a
learned scalar correction to $u_t$ (1{,}969 parameters). The GRU (hidden size
14), LSTM (12), and one-layer causal Transformer (width 16, four heads) have
3{,}664, 3{,}824, and 4{,}432 parameters. In Task~B the cyclic filter has 1{,}635
parameters and the GRU, LSTM, and Transformer have 7{,}312, 7{,}200, and 5{,}376.
All models train with per-step NLL on 10~s windows, Adam, batch size 32, and
gradient clipping at norm 1, for 6{,}000 updates in Task~A and 4{,}000 in
Task~B. The learning rate is chosen from $\{10^{-3},3\times10^{-3}\}$ by
validation NLL. Each setting uses five seeds.

\paragraph{Evaluation.}
Test windows start every 5~s. Their number is 739, 665, 558, 343, and 95 for
lengths 10, 30, 60, 120, and 240~s. All 18 test trials contribute except at
240~s. Only the three longest trials reach that length. We average seeds within each window
and then windows within each trial. Intervals resample whole test trials 10{,}000
times. VQF \citep{laidig2023vqf} runs offline on the raw 285.7~Hz data of each
window, in its 6D form for Task~A and its 9D form for Task~B.

\paragraph{Amendments.}
We fixed the protocol before any model was trained. In the first round no
learned model had a tilt estimate. All learned models had errors of
$25$--$39^\circ$ at 10~s, and the rotation with a learned angle collapsed to no
rotation. We kept this round as a negative result and added the gravity
direction and $u_t$ to the inputs of every learned model. The second round showed
the drift of learned angle corrections at 60~s and beyond. On validation trials
we traced the drift to a mean correction of $4.4\times10^{-4}$~rad per step.
We then added the input-set variant and ran it once on the test trials with the
same grid. We also added VQF as a reference. Task~B was specified before any of
its models were trained. Its first run read two feature columns in swapped
order. We discarded it and reran all Task~B models with corrected code. The
protocol specified two Transformer layers, but all reported runs use one layer.

\begin{table}[!htb]
\centering\small
\setlength{\tabcolsep}{4.5pt}
\caption{Task~A on held-out trials: heading error at the window end (degrees)
and test NLL (nats). Lower is better. Classical filters give point estimates
only.}
\label{tab:imu-a}
\begin{tabular}{lrrrrrcrrrrr}
\toprule
& \multicolumn{5}{c}{Error (deg)} && \multicolumn{5}{c}{NLL (nats)} \\
\cmidrule{2-6}\cmidrule{8-12}
Window length (s) & 10 & 30 & 60 & 120 & 240 && 10 & 30 & 60 & 120 & 240 \\
\midrule
Input-set rotation & 4.8 & 7.7 & 11.6 & 20.2 & 29.6 && 0.38 & 0.54 & 0.75 & 1.30 & 1.58 \\
Learned angle & 4.9 & 8.5 & 22.7 & 99.7 & 87.4 && 0.34 & 0.59 & 0.97 & 1.77 & 2.51 \\
Learned correction & 7.7 & 26.5 & 53.7 & 97.5 & 93.5 && 0.53 & 1.43 & 2.40 & 3.34 & 4.28 \\
GRU & 25.4 & 45.8 & 58.1 & 60.1 & 75.2 && 1.07 & 1.92 & 2.39 & 2.63 & 3.63 \\
LSTM & 23.0 & 43.5 & 53.9 & 53.6 & 73.9 && 1.00 & 1.86 & 2.37 & 2.66 & 4.13 \\
Transformer & 21.9 & 46.0 & 45.6 & 47.9 & 73.7 && 0.99 & 2.78 & 3.76 & 4.26 & 6.64 \\
\midrule
Integrated $u_t$ & 4.5 & 7.1 & 12.3 & 21.5 & 31.0 && \multicolumn{5}{c}{---} \\
Complementary filter & 6.9 & 9.1 & 13.8 & 22.5 & 31.4 && \multicolumn{5}{c}{---} \\
VQF (6D) & 4.2 & 6.0 & 7.7 & 6.8 & 20.7 && \multicolumn{5}{c}{---} \\
\bottomrule
\end{tabular}
\end{table}

\begin{table}[!htb]
\centering\small
\caption{Task~A paired differences, input-set rotation minus the comparison
(negative favors the input-set rotation), with 95\% trial-bootstrap intervals.
The best network is chosen per length and metric.}
\label{tab:imu-contrast}
\setlength{\tabcolsep}{4pt}
\begin{tabular}{rlrrrr}
\toprule
& & \multicolumn{2}{c}{vs.\ best network} & \multicolumn{2}{c}{Error (deg) vs.\ classical} \\
\cmidrule(lr){3-4}\cmidrule(lr){5-6}
Length (s) & Best network & Error (deg) & NLL (nats) & Integrated $u_t$ & VQF (6D) \\
\midrule
10 & Transformer & $-17.1$ [$-23.6$, $-11.3$] & $-0.62$ [$-0.83$, $-0.42$] & $+0.4$ [$+0.1$, $+0.6$] & $+0.6$ [$+0.4$, $+0.9$] \\
30 & LSTM & $-35.8$ [$-48.6$, $-23.1$] & $-1.32$ [$-1.69$, $-0.95$] & $+0.6$ [$+0.2$, $+1.1$] & $+1.7$ [$+1.2$, $+2.3$] \\
60 & Transf./LSTM & $-34.0$ [$-48.7$, $-19.3$] & $-1.61$ [$-2.11$, $-1.10$] & $-0.8$ [$-1.4$, $-0.2$] & $+3.8$ [$+2.6$, $+5.3$] \\
120 & Transf./GRU & $-27.7$ [$-44.2$, $-10.5$] & $-1.33$ [$-1.85$, $-0.78$] & $-1.3$ [$-2.3$, $-0.4$] & $+13.4$ [$+9.3$, $+17.2$] \\
240 & Transf./GRU & $-44.1$ [$-65.5$, $-31.4$] & $-2.05$ [$-2.67$, $-1.46$] & $-1.4$ [$-2.0$, $-0.6$] & $+9.0$ [$+6.9$, $+11.0$] \\
\bottomrule
\end{tabular}
\end{table}

\begin{table}[!htb]
\centering\small
\setlength{\tabcolsep}{4.5pt}
\caption{Task~B tracking error (degrees, averaged over valid steps) and test
NLL (nats) on magnetically disturbed (5) and undisturbed (13) test trials.
Classical filters give point estimates only.}
\label{tab:imu-b}
\begin{tabular}{lrrrrcrrrrcrr}
\toprule
& \multicolumn{4}{c}{Disturbed, error} && \multicolumn{4}{c}{Undisturbed, error} && \multicolumn{2}{c}{NLL, 60~s} \\
\cmidrule{2-5}\cmidrule{7-10}\cmidrule{12-13}
Window length (s) & 10 & 30 & 60 & 120 && 10 & 30 & 60 & 120 && Dist. & Undist. \\
\midrule
Cyclic filter (ours) & 9.0 & 8.9 & 9.5 & 8.9 && 5.9 & 6.0 & 6.1 & 6.1 && 2.64 & 2.57 \\
GRU & 9.9 & 11.0 & 12.6 & 11.5 && 5.5 & 5.9 & 6.2 & 6.2 && 2.67 & 1.99 \\
LSTM & 8.8 & 9.7 & 10.9 & 10.2 && 5.6 & 6.1 & 6.4 & 6.4 && 2.46 & 1.86 \\
Transformer & 9.2 & 11.9 & 15.8 & 18.0 && 5.0 & 6.7 & 8.7 & 10.9 && 3.20 & 3.40 \\
\midrule
Complementary filter & 10.8 & 9.5 & 9.2 & 8.4 && 7.2 & 7.2 & 7.3 & 7.3 && \multicolumn{2}{c}{---} \\
With rejection & 11.5 & 12.1 & 12.1 & 8.8 && 7.0 & 7.2 & 7.5 & 7.4 && \multicolumn{2}{c}{---} \\
VQF (9D) & 5.2 & 4.0 & 3.4 & 2.6 && 2.2 & 2.2 & 2.1 & 2.0 && \multicolumn{2}{c}{---} \\
\bottomrule
\end{tabular}
\end{table}

In Task~B the cyclic filter's belief update uses a learned von Mises pointer
likelihood with an outlier weight, and a learned diffusion of the shifted
belief. The criterion fixed in advance asked it to beat both the best network
and the tuned rejection filter on disturbed trials with intervals excluding
zero. At 60~s the differences are $-1.4^\circ$ $[-3.4,0.6]$ against the LSTM and
$-2.6^\circ$ $[-6.5,1.3]$ against the rejection filter, so the criterion is not
met. On undisturbed trials the filter is $1.2$--$1.4^\circ$ more accurate than
both complementary filters, with intervals excluding zero, and ties the GRU.
